% =====================================================================
%  22.5  The Engine Boundary: when Heavy-Light is the WRONG tool.
%  WORKED EXAMPLES ONLY -- NO \solvedproblem (in-section count stays 4).
%  Exports: sec:c16-boundary  def:c16-map  lst:c16-map
%  Every forward-cite here is PROSE + \chref only, NO code and NO re-link:
%    dynamic topology (link/cut)      -> Link-Cut / Euler-Tour, Ch.25 (c23)
%    STATIC path k-th / count         -> persistent four-root engine, Ch.21 (c19)
%    non-composable path updates      -> segment-tree beats / Li Chao, Ch.23 (c29)
%    OFFLINE path queries             -> Mo's on trees, Ch.27 (c21)
%    path COUNTING / distance stats   -> centroid decomposition, Ch.24 (c18), name-only
%    fully dynamic + path + subtree    -> top trees, name-only paragraph
%  Back-cites discharged promises: Ch.20 (c15) defers path updates HERE; Ch.21
%  (c19) defers the mutable path HERE -- this section returns the favour by
%  naming what those two chapters keep for themselves.
% =====================================================================
\section{The Engine Boundary: When Heavy-Light Is the Wrong Tool}\label{sec:c16-boundary}

\begin{previouslybox}[Four sections said what HLD is for; this one says what it is \emph{not} for]
Sections 22.1--22.4 built the chain split, the \(\mathit{pos}\) order, path queries, path updates, and the combined subtree-and-path structure.  Chapter~20 (\chref{c15}) handed its mutable-path limit to this chapter; Chapter~21 (\chref{c19}) handed its ``path that is \emph{updated}'' row here.  Those hand-offs are now discharged.  A boundary section exists because a powerful, easy-to-write tool invites over-use: heavy-light commits hard to a \emph{fixed} rooted topology, a \emph{static} decomposition, and \emph{composable} (monoid) updates.  Break any of those three assumptions and a different engine -- each already built elsewhere in this book -- strictly dominates.  This section draws the map so you reach for the right one on sight of a specification.
\end{previouslybox}

\begin{topicintro}
Three diagnostic questions place any tree-with-aggregates problem on the map. \emph{Is the topology fixed?} If edges/roots appear and disappear, HLD's chains shatter on every change and you need a dynamic-tree structure. \emph{Are the values mutable?} If nothing is ever updated and you only ask static path \(k\)-th / count questions, the persistent engine answers in fewer logs. \emph{Do the updates compose?} If an update is \code{chmin}/modulo/etc.\ rather than add/assign, a lazy tag cannot summarise it and beats-style analysis is required. Answering these three, plus whether the queries may be reordered (offline), picks the engine before you write a line.

\textbf{Prerequisites.} The whole of 22.1--22.4; Ch.~21 (\(\chref{c19}\)) for the static path contrast; Ch.~20 (\(\chref{c15}\)) for what a lazy tag can and cannot express; Ch.~11 (\(\chref{c11}\)) for the LCA that HLD reuses and that a dynamic-tree engine replaces.
\end{topicintro}

% ---------------------------------------------------------------
\subsection{The three axes that decide the engine}
% ---------------------------------------------------------------

\begin{definition}{The tree-aggregate engine map}{c16-map}
Every problem in this family is classified along three independent axes: \emph{topology} (fixed vs mutable -- links/cuts), \emph{values} (static vs mutable -- updates), and \emph{query} (a composable monoid fold vs a non-composable or order-statistic operation vs an offline batch).  Heavy-light owns the cell \(\{\text{fixed topology, mutable values, composable path/subtree fold, online}\}\).  Each neighbouring cell is a distinct, already-built engine.
\end{definition}

\begin{center}
\small
\renewcommand{\arraystretch}{1.28}
\setlength{\tabcolsep}{4pt}
\begin{tabular}{@{}p{0.20\linewidth} p{0.155\linewidth} p{0.14\linewidth} p{0.30\linewidth} >{\raggedright\arraybackslash}p{0.11\linewidth}@{}}
\toprule
\textbf{Symptom in the spec} & \textbf{Topology} & \textbf{Values} & \textbf{Right engine (and why HLD fails)} & \textbf{Chapter} \\
\midrule
add / assign a path, fold a path & fixed & mutable & \textbf{HLD + lazy} -- the home cell & this ch. \\
subtree \emph{and} path, mixed, online & fixed & mutable & \textbf{HLD order} (one \(\mathit{pos}\) tree) & \(\secref{c16-subpath}\) \\
link / cut edges, then query a path & \emph{mutable} & mutable & Link-Cut / Euler-Tour tree; HLD's DFS order dies each edit & Ch.~25 \(\chref{c23}\) \\
\(k\)-th / count on a path, \emph{never} updated & fixed & \emph{static} & persistent four roots: \(O(\log n)\) not \(O(\log^2 n)\); HLD pays an extra log & Ch.~21 \(\chref{c19}\) \\
path update is \code{chmin}/gcd/modulo & fixed & mutable & segment-tree beats; a lazy tag can't summarise a non-monoid op & Ch.~23 \(\chref{c29}\) \\
offline path queries, add/remove is cheap & fixed & either & Mo's on trees: \(O(n\sqrt n)\) beats \(O(\log^2 n)\) at high \(q\) & Ch.~27 \(\chref{c21}\) \\
count / sum over \emph{all} paths, no ranges & fixed & static & centroid decomposition; enumerate by centroid, not by chain & Ch.~24 \(\chref{c18}\) \\
\bottomrule
\end{tabular}
\end{center}

\begin{lstlisting}[language={}, basicstyle=\ttfamily\small, numbers=none, frame=single,
  caption={Decision ladder: read top to bottom, stop at the first line that matches the specification. Hands off everything HLD does not own.}, label={lst:c16-map}]
  Does the tree's EDGE SET or ROOT change between operations?
    |
    +-- YES -> NOT HLD.  Dynamic topology:
    |          * link/cut + path aggregates ........ Link-Cut Tree  (Ch.25, c23)
    |          * subtree-in-order + insert/delete .. Euler-Tour Tree (Ch.25, c23)
    |          (HLD re-runs two DFS passes per edit -> O(n) -> TLE.)
    |
    +-- NO (topology fixed).  Are the node/edge VALUES ever updated?
         |
         +-- NO (fully static).  Query is k-th / count-by-value on a path?
         |    +-- persistent roots + LCA (Ch.21, c19): O(log n), four roots.
         |       (HLD + merge-sort works but re-pays a log: O(log^3 n).)
         |
         +-- YES (mutable).  Is every update a COMPOSABLE monoid op
              (add, assign, multiply, max-with-identity)?
              |
              +-- YES, and queries are ONLINE, path and/or subtree:
              |    +>> HEAVY-LIGHT + LAZY  (this chapter) <-- your cell
              |         build O(n); op O(log^2 n); one pos-order tree.
              |
              +-- NO (chmin / modulo / gcd-update, "range floor-sum"):
              |    +-- segment-tree beats (Ch.23, c29): needs amortised
              |       potential, not a lazy tag. HLD can still supply the
              |       CHAIN ranges, but the per-range engine is a beats node.
              |
              +-- queries may be REORDERED (offline, many path queries):
                   +-- Mo's on trees (Ch.27, c21): O(n sqrt n), simpler add/
                      remove when the aggregate has no inverse.
\end{lstlisting}

% ---------------------------------------------------------------
\subsection{Worked boundary calls (read a spec, name the engine)}
% ---------------------------------------------------------------

\begin{comparisonbox}[Three specifications, three different answers -- none solved here]
\textbf{(A) ``n up to \(2\cdot10^5\); add an edge, delete an edge, then query the sum along the current path between two nodes, online.''}  The edge set is \emph{mutable}.  Every link/cut invalidates \(\mathit{dfs1}\)'s sizes and \(\mathit{dfs2}\)'s \(\mathit{pos}\), so an HLD rebuild costs \(\bigO{n}\) per edit and the whole thing degenerates.  This is a \term{Link-Cut Tree} (Ch.~25 \(\chref{c23}\)) problem: splay trees laid over preferred paths give \(\bigO{\log n}\) per link/cut/query with no static decomposition.

\textbf{(B) ``a fixed weighted tree; many queries `what is the \(k\)-th smallest node value on path \(u\)--\(v\)?'; no updates.''}  Topology fixed, values static, but the query is an order statistic.  HLD would decompose the path into \(O(\log n)\) chains and merge-sort-search each -- \(\bigO{\log^3 n}\).  Chapter~21 (\(\chref{c19}\)) plants one frequency root per tree node and answers the same query in \(\bigO{\log n}\) with the four-root descent.  With nothing to update, persistence's no-rebuild win is free; reach for it.

\textbf{(C) ``fixed tree; operation `set every value on path \(u\)--\(v\) to \(\min(\text{value},x)\)'; query the path sum.''}  A \code{chmin} is \emph{not} a monoid action -- it does not commute with a stored sum the way an add does, so a lazy tag has nothing to record.  This is \emph{segment-tree beats} (Ch.~23 \(\chref{c29}\); the term is indexed there).  The heavy-light skeleton is still useful -- it turns the path into chain ranges -- but the per-range engine must be a beats node, not Ch.~20's tag.
\end{comparisonbox}

\begin{observation}
HLD is a \emph{reduction}, not an aggregate engine: it converts ``a path/subtree on a tree'' into ``\(O(\log n)\) ranges on an array.''  Whatever array structure you then plug in -- plain sum tree, lazy tag (Ch.~20), beats node (Ch.~23) -- inherits the \(O(\log n)\) factor.  So the boundary is really about \emph{what breaks the reduction}: mutable topology breaks the \(\mathit{pos}\) order outright; static values make the reduction an unnecessary extra log; non-composable updates leave the reduction intact but demand a different leaf structure.
\end{observation}

\begin{notebox}[Top trees, in one honest paragraph (name-only)]
A \term{top tree} is the most general dynamic-tree structure: it clusters edges into a binary tree of \emph{clusters} by boundary vertices, so it supports link, cut, and aggregate over any path \emph{or} subtree in \(\bigO{\log n}\) -- strictly dominating both link-cut trees and HLD in full generality.  It is deliberately \emph{not} implemented in this book: the constants and case analysis are large, and in contest practice the specific subsumed query is almost always served by a lighter engine (HLD for fixed topology, LCT for links/cuts, centroid for path statistics).  Knowing it exists closes the map; you should reach past HLD to a top tree only when a problem needs dynamic topology \emph{and} subtree-and-path aggregates \emph{simultaneously} and defeats LCT.
\end{notebox}

\begin{stuckbox}[If you keep wanting HLD but the spec has a link/cut]
Stop and re-read the operation list.  A single ``add edge / remove edge / change the parent'' query means the decomposition is no longer amortised -- you would rebuild it every step.  That is the trigger to switch to Ch.~25 (\(\chref{c23}\)).  If instead the trouble is a strange update like ``modulo each value by \(p\),'', the topology is fine and the culprit is a non-monoid action: go to Ch.~23 (\(\chref{c29}\)).  Diagnosing which of the three axes (\(\defref{c16-map}\)) is violated is the whole skill this section teaches.
\end{stuckbox}

\begin{recallbox}[22.5 in one breath]
\emph{Heavy-light owns exactly one cell of the map: fixed topology, mutable values, composable online path/subtree folds.  Mutable topology \(\to\) link-cut / Euler-tour (Ch.~25).  Fully static path order-statistics \(\to\) persistence (Ch.~21).  Non-composable updates \(\to\) beats (Ch.~23).  Offline \(\to\) Mo on trees (Ch.~27).  Whole-tree path statistics \(\to\) centroid (Ch.~24).  Top trees subsume the dynamic cases but are not implemented.}
\selftest{\item A problem fixes the tree but asks ``path \(k\)-th smallest'' with no updates at all -- name the engine and the reason you do \emph{not} use HLD. \emph{Persistent four-root tree (Ch.~21, \(\chref{c19}\)); HLD would answer correctly but re-pays a chain-log (\(\bigO{\log^3 n}\) vs \(\bigO{\log n}\)) and needs no persistence, so with nothing mutable the reduction is pure overhead.}}
\end{recallbox}
